% Generated draft: requires mathematical and source review.
\documentclass[11pt]{article}
\usepackage[margin=25mm]{geometry}
\usepackage{fontspec}
\setmainfont{Latin Modern Roman}
\usepackage{amsmath,amssymb,mathtools,unicode-math}
\setmathfont{Latin Modern Math}
\usepackage{enumitem,needspace,xurl,hyperref}
\setlength{\parindent}{0pt}
\setlength{\parskip}{5pt}
\setlength{\emergencystretch}{3em}
\newcommand{\N}{\mathbb{N}}
\newcommand{\Z}{\mathbb{Z}}
\newcommand{\Q}{\mathbb{Q}}
\newcommand{\R}{\mathbb{R}}
\newcommand{\C}{\mathbb{C}}
\newcommand{\F}{\mathbb{F}}
\newcommand{\A}{\mathbb{A}}
\newcommand{\PP}{\mathbb{P}}
\newcommand{\E}{\mathbb{E}}
\newcommand{\eps}{\varepsilon}
\newcommand{\dd}{\,\mathrm{d}}
\begin{document}
\section*{UnsolvedMath ID: 20003120}
\textbf{Source problem code:} catalogue entry
\subsection{Definitions and mathematical statement}
For a Riemann surface $X$, a conformal metric has local form $\rho(z)|\dd z|$, with lengths $\ell_\rho(\gamma)=\int_\gamma\rho|\dd z|$ and area $\operatorname{Area}(\rho)=\int_X\rho^2\,\dd x\dd y$. The source asks whether there is a least-area metric satisfying
\[
\inf_{\gamma\in\Gamma_X}\ell_\rho(\gamma)\geq1,
\]
and whether Thurston's metric is related to this problem. The aid interprets $\Gamma_X$ as all homotopically nontrivial closed curves, including puncture-peripheral curves in the punctured case. Smooth versus measurable or singular admissible metrics, compactness, and renormalization are not specified by the source and affect attainment.

For a closed surface of genus at least one define the positive finite infimal value
\[
A_{\min}(X)=\inf_{\rho:\ \inf_{\gamma\in\Gamma_X}\ell_\rho(\gamma)\geq1}\operatorname{Area}(\rho).
\]
With the matching admissible class, generalized extremal length is
\[
\operatorname{EL}_X(\Gamma_X)=\sup_\rho
\frac{(\inf_{\gamma\in\Gamma_X}\ell_\rho(\gamma))^2}{\operatorname{Area}(\rho)},
\qquad A_{\min}(X)=\operatorname{EL}_X(\Gamma_X)^{-1}.
\]
For two marked conformal structures, the supplied quasiconformal comparison gives
\[
|\log A_{\min}(Y)-\log A_{\min}(X)|\leq2d_T(X,Y).
\]
For closed genus at least two, on a fixed $\epsilon$-thick part the aid invokes Choi--Rafi to conclude
\[
|\log A_{\min}(Y)-\log A_{\min}(X)|\leq
2\max\{d_{\mathrm{Th}}(X,Y),d_{\mathrm{Th}}(Y,X)\}+2C(S,\epsilon).
\]
The topology/thickness constant is not supplied numerically. These inequalities concern the infimum and do not assume or prove that it is attained.

If loops around a puncture are constrained to length at least one, a coordinate annulus $\epsilon<|z|<r_0$ has
\[
\operatorname{Area}_\rho(\epsilon<|z|<r_0)\geq\frac1{2\pi}\log\frac{r_0}{\epsilon}.
\]
The model cylinder density $\rho=1/(2\pi|z|)$ saturates the coefficient. Thus ordinary area necessarily diverges logarithmically, with coefficient $1/(2\pi)$ per constrained puncture, matching the aid's Wolf--Zwiebach renormalization convention. A finite-part minimization requires an explicit subtraction and cutoff-coordinate convention; it is a different target from finite ordinary area. Excluding peripheral loops would also be a different constraint.

For a torus $X_\tau=\C/(\Z+\tau\Z)$, $\operatorname{Im}\tau>0$, the aid reports a unique normalized flat minimizer. The coordinate-invariant value is
\[
A_{\min}(X_\tau)=\frac{\operatorname{Im}\tau}{s(\tau)^2},\qquad
s(\tau)=\min_{(a,b)\in\Z^2\setminus\{0\}}|a+b\tau|.
\]
In a standard modular fundamental-domain normalization $s(\tau)=1$, this is the reported $A_{\min}(X_\tau)=\operatorname{Im}\tau$. The missing shortest-period normalization is essential: the latter formula is not invariant for an arbitrary representative $\tau$. Uniqueness is in the appropriate conformal-density sense, up to the metric-class convention. The genus-zero empty-constraint case and general attainment in the intended surface/regularity class are not resolved by these special results.
\subsection{Short English statement}
The supplied result controls the minimal-area value through generalized extremal length and Teichm\"uller distance, with a thick-part comparison to the symmetric Thurston distance. This does not prove a minimizer exists. Peripheral constraints at punctures force logarithmically infinite ordinary area, while the torus case has a normalized flat minimizer. The torus formula needs a shortest-period normalization, and punctures need a separate renormalized problem.
\subsection{Sources}
\begin{itemize}
\item Local registry: permanent catalogue entry, exact source code catalogue entry; \texttt{aim-topology-notes.json}, \texttt{notes[207]}.
\item Weixu Su's AIM list \emph{Problems on Thurston metric}, Problem 5.2 attributed to D. Dumas; workshop catalog title \emph{Lipschitz metric on Teichmueller space}, \url{https://aimath.org/pastworkshops/lipschitzteichproblems.pdf}. The source refers to work of Wolf and Zwiebach and applications to string theory, without full publication details.
\item Supplied attribution: Choi--Rafi thick-part metric comparison, without title, theorem number, or complete citation. The original remarks and separate literature field are empty.
\item Supplied research attempt 1, registry updated 14 August 2026: partial result, open-at-run label, low novelty confidence, machine-generated and not peer reviewed. It supplies the value inequalities, puncture coefficient, and torus case. Verification reports that the official formula repairs visible extraction fragments, but ends before reproducing its checked display.
\end{itemize}
\end{document}
